\DocumentMetadata{
  pdfversion=2.0,
  pdfstandard=ua-2,
  lang=en-US,
  tagging=on,
  tagging-setup={math/setup=mathml-SE,tabsorder=structure}
}
\documentclass[11pt,letterpaper]{article}
\usepackage[margin=0.68in,includefoot]{geometry}
\usepackage{newcomputermodern}
\usepackage{amsmath,amscd,mathtools}
\usepackage{tikz,adjustbox}
\providecommand{\square}{\mdlgwhtsquare}
\usepackage[hidelinks]{hyperref}
\hypersetup{
  pdftitle={Complex Analysis PhD Exam, September 2011},
  pdfauthor={University of Florida Department of Mathematics},
  pdfsubject={Graduate mathematics examination},
  pdfdisplaydoctitle=true,
  bookmarksopen=true
}
\setcounter{secnumdepth}{0}
\setlength{\parindent}{0pt}
\setlength{\parskip}{0.28em}
\setlength{\emergencystretch}{3em}
\setlength{\leftmargini}{2.15em}
\setlength{\leftmarginii}{2.65em}
\setlength{\leftmarginiii}{3em}
\pagestyle{empty}
\raggedbottom
\allowdisplaybreaks
\NewTaggingSocketPlug{block/list/label}{examproblem}{%
  \tagstructbegin{tag=Lbl}\tagstructend%
  \tagstructbegin{tag=\LBody}%
  \tagstructbegin{tag=\ProblemHeadingTag,title-o={\ProblemHeadingText}}%
  \tagmcbegin{tag=\ProblemHeadingTag,actualtext={\ProblemHeadingText}}%
  #2%
  \tagmcend%
  \tagstructend%
}
\newenvironment{examproblems}[2]{%
  \def\ProblemHeadingTag{#2}%
  \AssignTaggingSocketPlug{block/list/label}{examproblem}%
  \begin{enumerate}%
  \setcounter{enumi}{#1}%
  \renewcommand{\labelenumi}{\textbf{\theenumi.}}%
  \setlength{\topsep}{0.35em}%
  \setlength{\partopsep}{0pt}%
  \setlength{\itemsep}{0.48em}%
  \setlength{\parsep}{0.18em}%
}{\end{enumerate}}
\newenvironment{examsubparts}{%
  \AssignTaggingSocketPlug{block/list/label}{default}%
  \begin{enumerate}%
  \setlength{\topsep}{0.18em}%
  \setlength{\partopsep}{0pt}%
  \setlength{\itemsep}{0.16em}%
  \setlength{\parsep}{0.1em}%
}{\end{enumerate}}
\newenvironment{examinstructions}{%
  \par\begingroup\itshape
}{%
  \par\endgroup\vspace{0.18em}
}
\makeatletter
\renewcommand\subsection{\@startsection{subsection}{2}{0pt}%
  {-1.25ex plus -.25ex minus -.1ex}%
  {0.55ex plus .1ex}%
  {\normalfont\large\bfseries}}
\renewcommand{\@maketitle}{%
  \newpage\null\vskip -1em%
  {\footnotesize\bfseries
    \href{https://www.ufl.edu/}{University of Florida}, \href{https://math.ufl.edu/}{Department of Mathematics}\par}%
  \vskip -0.5em%
  \tagstructbegin{tag=H1,title={Complex Analysis PhD Exam, September 2011}}%
  \tagpdfparaOff%
  \tagmcbegin{tag=H1}%
    {\LARGE\bfseries\href{https://gma.math.ufl.edu/exams/complex-analysis-phd/}{Complex Analysis PhD Exam}, September 2011\par}%
  \tagmcend%
  \tagpdfparaOn%
  \tagstructend%
  \vskip 0.25em%
  \tagmcbegin{artifact}%
    \hrule height 1pt%
  \tagmcend%
  \vskip 1em%
}
\makeatother
\title{Complex Analysis PhD Exam, September 2011}
\author{}
\date{}
\begin{document}
\maketitle
\thispagestyle{empty}
\begin{examinstructions}
Do all 9 problems. The symbol~\(U\) always denotes the unit disc: \(|z|<1\); \(\mathbb{C}\), the complex plane; \(\Omega\), an open region in~\(\mathbb{C}\); and \(\Re z\), the real part of~\(z\).
\end{examinstructions}
\begin{examproblems}{0}{H2}
\def\ProblemHeadingText{Problem 1.}
\item
Evaluate the integral \(\displaystyle \int_0^\infty \frac{(\log x)^3}{1+x^2}\,dx\).
\def\ProblemHeadingText{Problem 2.}
\item
Let \(f_n\), \(n=1,2,3,\ldots\), be a sequence of one-to-one analytic functions on~\(\Omega\). If \(f_n\) converges uniformly to~\(f\) on any compact subset of~\(\Omega\), prove that~\(f\) is either constant or one-to-one. Show by examples that both conclusions can occur.
\def\ProblemHeadingText{Problem 3.}
\item
Let~\(F\) be a family of analytic functions~\(f\) on~\(U\) such that \(\Re f>0\) and \(f(0)=1\). Prove that~\(F\) is a normal family. Can the condition \(f(0)=1\) be omitted? Justify your answer.
\def\ProblemHeadingText{Problem 4.}
\item
Suppose~\(f\) is a conformal mapping of~\(U\) onto a square with center at~\(0\) and \(f(0)=0\). Prove that \(f(iz)=if(z)\). If \(f(z)=\sum_0^\infty c_nz^n\), prove that \(c_n=0\) unless \(n-1\) is a multiple of~\(4\).
\def\ProblemHeadingText{Problem 5.}
\item
Let~\(R\) be a rational function such that \(|R(z)|=1\) if \(|z|=1\). Prove that~\(R\) is of the form \(R(z)=cz^m\prod_{n=1}^k\frac{z-a_n}{1-z\overline{a_n}}\), where~\(c\) is a constant with \(|c|=1\), \(m\) is an integer and \(a_1,a_2,\ldots,a_k\) are complex numbers such that \(a_n\ne 0\) and \(|a_n|\ne 1\).
\def\ProblemHeadingText{Problem 6.}
\item
Let~\(f\) be analytic in~\(\Omega\). If~\(f\) has no zeros in~\(\Omega\), prove that \(\log|f|\) is harmonic in~\(\Omega\).
\def\ProblemHeadingText{Problem 7.}
\item
Suppose that \(f:\mathbb{C}\to\mathbb{C}\) is a continuous function and~\(f\) is analytic everywhere except possibly on \([-1,1]\). Prove that~\(f\) is an entire function.
\def\ProblemHeadingText{Problem 8.}
\item
Suppose~\(f\) and~\(g\) are both analytic in~\(U\), and neither of them has a zero in~\(U\). If \(\frac{f'}{f}(1/n)=\frac{g'}{g}(1/n)\) for \(n=2,3,4,\ldots\), find a simple relation between~\(f\) and~\(g\).
\def\ProblemHeadingText{Problem 9.}
\item
Construct a conformal map \(w=f(z)\) which maps~\(U\) onto the angular region \(A:|\arg(w)|<\alpha\) and \(f(0)=1\), where \(0<\alpha<\pi\).
\end{examproblems}
\end{document}
